% !TEX root = master.tex
% !TEX program = lualatex
\lecture{3}{10-03-2026}{Integration and Cauchy theorem}{}

\chapter{Complex integration}\label{chap:complex_integration}
\section{Line Integrals in $\mathbb{R}^2$ (Recall)}

Let $F : \mathbb{R}^2 \rightarrow \mathbb{R}^2$ be a continuous vector field
\[
F(x_1,x_2) = (F_1(x_1,x_2),F_2(x_1,x_2))
\]

Let $\gamma : [a,b] \rightarrow \mathbb{R}^2$ be a $C^1$ curve with image $\Gamma$.

The line integral of $F$ along $\gamma$ is defined as

\[
\int_\Gamma F \cdot d\ell
=
\int_a^b F(\gamma(t)) \cdot \gamma'(t)\,dt
\]

which explicitly gives

\[
\int_a^b \left[
F_1(\gamma(t))\gamma_1'(t) +
F_2(\gamma(t))\gamma_2'(t)
\right]dt
\]

\subsection*{Remarks}

\begin{itemize}
\item The line integral is independent of the parametrization.
\item Reversing the orientation changes the sign.
\item A curve is called
\begin{itemize}
\item \textbf{simple} if it has no self-intersections
\item \textbf{regular} if $\gamma'(t) \neq 0$
\end{itemize}
\end{itemize}

If $\Gamma$ is closed and $\nabla \times F = 0$ in the interior of $\Gamma$ then

\[
\int_\Gamma F \cdot d\ell = 0
\]

(Stokes' theorem).

\section{Complex Line Integrals}

Let $\gamma : [a,b] \to \mathbb{C}$ be a regular curve parametrizing $\Gamma \subset \mathbb{C}$.

Let $f : \Gamma \to \mathbb{C}$ be continuous.

\begin{definition}
The line integral of $f$ along $\Gamma$ is

\[
\int_\Gamma f(z)\,dz
=
\int_a^b f(\gamma(t))\gamma'(t)\,dt
\]
\end{definition}

\subsection*{Orientation}

If $-\Gamma$ denotes the curve with opposite orientation then

\[
\int_{-\Gamma} f(z)\,dz
=
-\int_\Gamma f(z)\,dz
\]

\subsection*{Remark}

The value of the line integral does not depend on the parametrization as long as the orientation is preserved.

\section{Examples}

\subsection{Integral of $z^2$ on the unit circle}

Let $\Gamma$ be the unit circle parametrized by

\[
\gamma(t) = e^{it}, \quad t \in [0,2\pi]
\]

Then

\[
\gamma'(t) = i e^{it}
\]

We compute

\[
\int_\Gamma z^2\,dz
=
\int_0^{2\pi} (e^{it})^2 (i e^{it})\,dt
\]

\[
=
i \int_0^{2\pi} e^{3it}dt
\]

\[
=
i \left[\frac{e^{3it}}{3i}\right]_0^{2\pi} = 0
\]

\subsection{Integral of $\frac{1}{z}$}

\[
\int_\Gamma \frac{1}{z}dz
=
\int_0^{2\pi}
\frac{1}{e^{it}}(i e^{it})dt
=
i\int_0^{2\pi}dt
=
2\pi i
\]

\subsection{Integral of $\frac{1}{z^2}$}

\[
\int_\Gamma \frac{1}{z^2}dz
=
\int_0^{2\pi}
\frac{1}{e^{2it}}(i e^{it})dt
=
i\int_0^{2\pi} e^{-it}dt
=
0
\]

\section{Cauchy's Theorem}

\begin{theorem}
Let $U \subset \mathbb{C}$ be open and simply connected.

Let $\Gamma$ be a closed piecewise regular curve contained in $U$.

If $f : U \to \mathbb{C}$ is holomorphic then

\[
\int_\Gamma f(z)\,dz = 0
\]
\end{theorem}

\textbf{Simply connected:} the domain has no holes.

Examples:

\[
\mathbb{C}, \quad \{z \in \mathbb{C} : |z| < 1\}
\]

are simply connected, while

\[
\mathbb{C}^* = \mathbb{C}\setminus\{0\}
\]

is not.

\textbf{Piecewise regular:} $\Gamma$ might have finitely many curves, e.g. square. 
\subsection{Idea of the Proof}

Write

\[
f(z) = u(x,y) + i v(x,y)
\]

and parametrize

\[
\gamma(t) = \alpha(t) + i\beta(t)
\]

Then

\[
\int_\Gamma f(z)\,dz
=
\int_a^b
\left[u(\gamma(t))+iv(\gamma(t))\right]
(\alpha'(t)+i\beta'(t))dt
\]

Separating real and imaginary parts gives two line integrals in $\mathbb{R}^2$.

Using Green's theorem and the Cauchy--Riemann equations

\[
\frac{\partial u}{\partial x} = \frac{\partial v}{\partial y},
\qquad
\frac{\partial u}{\partial y} = -\frac{\partial v}{\partial x}
\]

both integrals vanish, hence

\[
\int_\Gamma f(z)\,dz = 0.
\]

\subsection{Extension of Cauchy's Theorem}

\begin{corollary}
Let $\Gamma_0$ and $\Gamma_1$ be closed simple regular curves with

\[
\Gamma_1 \subset \text{int}(\Gamma_0)
\]

If $f$ is holomorphic on the region between them then
  \[
\int_{\Gamma_0} f(z)\,dz
=
\int_{\Gamma_1} f(z)\,dz
\]


\end{corollary}
This allows computation of complicated integrals using simpler curves.

\section{Cauchy Integral Formula}

\begin{theorem}
Let $D \subset \mathbb{C}$ be open and $f : D \to \mathbb{C}$ holomorphic.

Let $\Gamma$ be a closed simple curve inside $D$ oriented clockwise and let $z_0$ lie inside $\Gamma$.

Then

\[
f(z_0)
=
\frac{1}{2\pi i}
\int_\Gamma
\frac{f(z)}{z-z_0}dz
\]
\end{theorem}
This formula is particularly useful to compute integrals of the form

\[
\int_\Gamma \frac{f(z)}{z-z_0}\,dz.
\]
\\
\textbf{Remark:} if the curves is oriented counterclockwise we change the sign as usual.

If more than one $z_0$, we can use this formula if only one $z_0$ is inside $\Gamma$, if yes use this formula by using this $z_0$. 

I.e: 
\[
I = \int_{\gamma}^{} \frac{e^z}{(z-1)(z^2 + 4)} dx \quad \gamma = e^{it} \quad t \in [\ 0,2 \pi ]\
\]
Here $z_0 = 2i, -2i$ is out of the unit circle, so only do for $z_0 = 1$ \[
  f(z) = \frac{e^z}{z^2 + 4}
\]
\subsection*{Idea of the proof}

Consider a small circle around $z_0$

\[
\gamma_r(t) = z_0 + re^{it}
\]

Then

\[
\int_\Gamma \frac{f(z)}{z-z_0}dz
=
\int_{\gamma_r} \frac{f(z)}{z-z_0}dz
\]

which becomes

\[
\int_0^{2\pi} f(z_0 + re^{it}) i dt
\]

As $r \to 0$, continuity of $f$ gives

\[
2\pi i f(z_0)
\]

which proves the formula.

